\section{研究问题、实验范围与阅读地图}

\subsection{本文回答的四个研究问题}

这版稿子不再按“family 逐段展开”的方式组织，而是直接围绕四个 research questions 展开：

\begin{table}[H]
\centering
\small
\begin{tabularx}{\textwidth}{p{2.1cm} p{3.9cm} Y}
\toprule
RQ & 核心问题 & 本文使用的主要证据 \\
\midrule
RQ1 & Overall performance 到底落在什么版图上 & Instruments 主表、主曲线、baseline 与 hint-enabled frontier 对照 \\
RQ2 & Prefix hint 在训练里是否真的带来有效信号；不同 hint strategy 怎么比较 & rule-only vs hint-conditioned 主线、prefix-hint 2x2、max1 budget 敏感性 \\
RQ3 & 训练任务范围与 loss 设计各自会带来什么变化 & 3-task / sid-only / two-task task-scope ablation，以及 fixed vs fixed+CE vs full-sequence SFT+GRPO 状态表 \\
RQ4 & CE loss 系数如何影响后训练表现 & Instruments 与 Arts 上的 \texttt{0.001 / 0.005 / 0.01} fixed CE 对照 \\
\bottomrule
\end{tabularx}
\caption{本文的 RQ 结构。}
\end{table}

\subsection{共享实验背景}

\begin{itemize}
\item 主数据集：\texttt{Instruments-grec}
\item RL 起点：\texttt{Instruments-grec-sft-qwen4B-4-256-dsz0/checkpoint-495}
\item 主 launcher 目录：\path{hope/Qwen2_5-3B-Isntruct-qwen4B-4-256-MIMIGenRec-grec/}
\item RQ4 额外引入的第二数据集：\texttt{Arts-grec}，只用来分析 fixed CE 系数的跨数据集稳定性
\item 主读法：默认先看 best \texttt{NDCG@10}，再看对应的 \texttt{HR@50}、必要时再看整条 epoch 曲线
\end{itemize}

\subsection{结果来源与默认命名}

本文的表格与曲线统一读取本地同步回来的 checkpoint 指标文件，主要来自各实验目录下的 \path{checkpoint-*/metrics.json}。为避免在 RQ2/RQ3 里反复解释名称，正文统一采用下面的缩写：

\begin{itemize}
\item \texttt{dynamic gather-fix}：当前默认的 canonical adaptive hinting 实现
\item \texttt{fixed taskfix}：当前默认的 fixed task-aware prefix hint，实现上对应“ours”主线
\item \texttt{fixed taskfix sid-only}：fixed first-token hint
\item \texttt{dynamic sid-only}：adaptive first-token hint
\item old \texttt{fixed mixed-single}：带 bug caveat 的历史参考线，正文只在必要处提醒，不再把它当 clean evidence
\end{itemize}

\subsection{正文与附录的边界}

\begin{itemize}
\item 正文只保留直接服务 RQ1--RQ4 的主证据。
\item old fixed bug、single-hint mixed、single-hint mixed + CE、hint-token SFT，以及 Games 的 cross-dataset readout，统一放到 appendix。
\item 因此正文的目标不是“把所有 exploratory line 全都讲一遍”，而是把 prefix hint 这条主故事压成一条可答辩的 RQ 链。 
\end{itemize}
